\begin{abstract}
Treat the realised final layer as what it is: $n$ independent Gaussian random-feature measurements of one
state, the law $\lambda$ of the penultimate activations. Three things follow.
\begin{enumerate}[label=(\arabic*),leftmargin=1.6em,itemsep=1pt,topsep=1pt]
\item \emph{Exact RMS as a kernel distance} (Thm.~\ref{thm:mmd}). Let $\lambda'$ be any signed surrogate state
independent of $W_L$; every estimator of stages~14--22 is the readout of one. Then
$\E[\mathrm{RMS}^2\mid\mathcal F]=\mathrm{MMD}^2_{k_1}(\lambda,\lambda')$, the maximum mean discrepancy for the degree-$1$
arc-cosine kernel. The realised RMS concentrates around it at rate $n^{-1/2}$. The error ledger of stage~18 is
exactly the Mercer expansion of this distance.
\item \emph{Transverse analyticity} (Thm.~\ref{thm:analytic}). The arc-cosine kernel is singular only on
coincident directions. Expanded around the collective angle, its Taylor series in the transverse overlap
$\rho_\perp$ of two inputs has nonnegative coefficients. The tail beyond order $J$ is bounded by $\pi\,|\rho_\perp|^{J+1}$.
So for the RMS, the kink is averaged over the random final weights, and the obstruction of stage~20, which
concerns a single neuron's moments, does not apply.
\item \emph{What a real network demands} (\S\ref{sec:real}). We measured one realised network ($n=512,1024$;
$L=16$). The transverse field is strongly anisotropic: its effective dimension is $n_{\rm eff}\approx80$ at $n=1024$,
so $|\rho_\perp|\approx0.11$, not $n^{-1/2}=0.03$. The per-neuron standardised skewness is $\approx0.12$ RMS. The absolute
$\kappa_3$ is $O(1/n)$, as before, but the collapsed transverse variance ($\sigma_k^3\approx0.06$) amplifies it. This
corrects earlier size claims. The $\kappa_3$ term moves each mean by $\approx10^{-3}$ RMS, sixteen times the target,
so the history readout must be accurate to about $6\%$. The $10\kappa_3^2$ term of $m_6$, which the audit flagged,
and the incoherent $\kappa_4$ are each at the target scale and must be kept.
\end{enumerate}
We also prove that each random layer adds $1/n$ to the inverse participation ratio
(Prop.~\ref{prop:pr}). Measured, $n/n_{\rm eff}$ grows linearly in depth with slope $0.85$. So the per-neuron expansion
parameter is the depth-to-width ratio $\sqrt{0.85L/n}\approx0.12$, not $n^{-1/2}$. Two adaptive checks follow from the
probe block: an order selector from $n_{\rm eff}$, and a skewness calibration test with a slope standard error of
about $0.02$ (\S\ref{sec:adapt}).
\end{abstract}
\maketitle

\section{The final layer measures a state}
Conditional on $\mathcal F=\sigma(W_1,\dots,W_{L-1})$, the columns $w_k\sim N(0,s^2I)$ are i.i.d. and independent of the penultimate
law $\lambda$. Put $y(w)=\int(w\cdot u)_+\,d\lambda(u)$, so $y_k=y(w_k)$. For a finite signed measure $\lambda'$ independent of
$W_L$, put $\hat y(w)=\int(w\cdot u)_+\,d\lambda'(u)$.

\begin{theorem}[Exact RMS as an arc-cosine MMD]\label{thm:mmd}
Let $k_1(u,u')=\E_w[(w\cdot u)_+(w\cdot u')_+]=\frac{s^2}{2\pi}\|u\|\|u'\|\,\kappa(\cos\theta)$, where
$\kappa(t)=\sqrt{1-t^2}+(\pi-\arccos t)\,t$. Then
\[
\E\big[\mathrm{RMS}(\hat y-y)^2\,\big|\,\mathcal F\big]=\mathrm{MMD}^2_{k_1}(\lambda,\lambda')=\iint k_1\,d(\lambda-\lambda')^{\otimes2},
\]
and $\mathrm{Var}(\mathrm{RMS}^2\mid\mathcal F)=\mathrm{Var}_w\big((\hat y-y)^2\big)/n$. Moreover:
\begin{enumerate}[label=(\alph*),leftmargin=1.6em,itemsep=1pt,topsep=1pt]
\item \emph{Mercer equals ledger.}
\[
\mathrm{MMD}^2=s^2\sum_{r\ge0}\frac{c_r^2}{r!}\,\big\|M_r(\lambda)-M_r(\lambda')\big\|_F^2,\qquad M_r(\nu)=\int\|u\|^{1-r}u^{\otimes r}d\nu ,
\]
with $c_r$ the Hermite coefficients of $\relu$. This is stage~18's Wiener-chaos ledger in the final weights.
\item \emph{All our estimators are signed-state readouts.} The Gaussian closure is the readout of $N(\mu,C)$.
The cumulant corrections are readouts of the multivariate Edgeworth signed measure, whose $\kappa_r$ tensors
are the hidden cumulants. The collective form is the readout of a mixture of these.
\end{enumerate}
\end{theorem}
\begin{proof}
$\hat y-y=\int(w\cdot u)_+\,d(\lambda'-\lambda)$. Squaring, taking $\E_w$ and using Fubini gives the MMD. The $n$ columns are
i.i.d., which gives the variance. The arc-cosine formula is the classical Gaussian computation.

For (a), the Mehler expansion $(w\cdot u)_+=s\|u\|\sum_rc_r\mathrm{He}_r(w\cdot\hat u/s)/r!$ and Hermite orthogonality give
$k_1=s^2\|u\|\|u'\|\sum_r(c_r^2/r!)(\hat u\cdot\hat u')^r$. Integrate against $(\lambda-\lambda')^{\otimes2}$.

For (b), the one-dimensional marginal of a multivariate Edgeworth measure along $w$ is the one-dimensional
Edgeworth measure with cumulants $\langle\kappa_r,w^{\otimes r}\rangle$. Its $\relu$ readout is stage~17's $F$.
\end{proof}

In NCG terms, the kernel $k_1$ is the GNS inner product of the ``final-layer state''. The mean embedding
$\xi_\lambda=\int\phi_u\,d\lambda(u)$, with $\phi_u=(w\mapsto(w\cdot u)_+)$, is a vector in $L^2(\gamma_s)$. The final means are its overlaps
with $n$ random Gaussian measurement vectors. As with classical shadows, random measurements estimate a
distance between states: the RMS is the GNS distance $\|\xi_\lambda-\xi_{\lambda'}\|$ up to $n^{-1/2}$ fluctuations. The
design question becomes: find a computable signed state $\lambda'$ at the penultimate layer within GNS
distance $6\times10^{-5}$ of the true one.

\section{Transverse analyticity: why the RMS escapes the moment obstruction}
Write each penultimate vector as $u=R(\cos\beta\,\hat\mu+\sin\beta\,\hat v)$, with $\hat v\perp\hat\mu$ (stage~18), and for a pair put
$\rho_\perp=\hat v\cdot\hat v'$. Then
\[
\cos\theta=t_0+h,\qquad t_0=\cos\beta\cos\beta',\qquad h=\sin\beta\sin\beta'\,\rho_\perp .
\]

\begin{lemma}[Positivity]\label{lem:pos}
$\kappa(t)=\sum_{m\ge0}b_mt^m$ with $b_m=2\pi c_m^2/m!\ge0$ and $\kappa(1)=\pi$. For $t_0\in[0,1)$, every Taylor coefficient
$\kappa^{(j)}(t_0)/j!=\sum_mb_m\binom mjt_0^{m-j}$ is nonnegative, and $\sum_j\kappa^{(j)}(t_0)(1-t_0)^j/j!=\kappa(1)$.
\end{lemma}
\begin{proof}
The expansion is Mehler's, as in Theorem~\ref{thm:mmd}(a), normalised so that $k_1(u,u)=s^2\|u\|^2/2$. The series has
nonnegative coefficients and converges at $t=1$, so Taylor re-expansion at $t_0$ is valid on $|h|\le1-t_0$.
\end{proof}

\begin{theorem}[Transverse expansion with geometric tail]\label{thm:analytic}
For every pair, $|h|\le1-t_0$. Splitting $\kappa(t_0+h)=\sum_{j\le J}\kappa^{(j)}(t_0)h^j/j!+\mathcal T_J$, we have
\[
0\le\mathcal T_J\ \text{in the PD sense},\qquad|\mathcal T_J|\le\pi\,|\rho_\perp|^{J+1}.
\]
Consequently
$\mathrm{MMD}^2_{k_1}(\lambda,\lambda')=\sum_{j\le J}\sum_m\frac{s^2b_m}{2\pi}\binom mj\|\Delta_{jm}\|^2+E_{>J}$, where
\[
\Delta_{jm}=\int R^{1-m}P^{m-j}u_\perp^{\otimes j}\,d(\lambda-\lambda'),\qquad
0\le E_{>J}\le\frac{s^2}2\iint RR'|\rho_\perp|^{J+1}\,d|\lambda-\lambda'|^{\otimes2}.
\]
\end{theorem}
\begin{proof}
$1-t_0-|h|\ge1-\cos(\beta-\beta')\ge0$, so $|h|\le1-t_0$. Also $|h|/(1-t_0)\le|\rho_\perp|$, because $1-\cos\beta\cos\beta'\ge\sin\beta\sin\beta'$. By
Lemma~\ref{lem:pos} the tail is a nonnegative combination of powers of $h$, bounded by
$(|h|/(1-t_0))^{J+1}\sum_j\kappa^{(j)}(t_0)(1-t_0)^j/j!$. The low part is a finite sum of products of feature maps. Indeed
$R\cos^{m-j}\beta\sin^j\beta\,\hat v^{\otimes j}=R^{1-m}P^{m-j}u_\perp^{\otimes j}$, which gives the $\Delta_{jm}$.
\end{proof}

\noindent\textbf{What this means.}
\begin{itemize}[leftmargin=1.2em,itemsep=1pt,topsep=1pt]
\item The arc-cosine kernel is singular only where two hidden directions coincide, which has probability
zero for pairs of independent inputs. On the support of such pairs it is analytic, with expansion parameter
$|\rho_\perp|$.
\item The per-neuron obstruction of stage~20 concerns one fixed $w$, which probes the kink. The RMS is an
average over $n$ independent $w$, which smooths the kink into $k_1$.
\item The order-$j$ terms vanish exactly when $\lambda'$ matches the transverse moment tensors of order $j$
\emph{jointly with} the collective coordinates (all $m$). That is stage~18's $B_j(a)$ matching.
\item A certificate needs two things: bounds on the mismatches $\|\Delta_{jm}\|$ for $j\le J$, and a pair-overlap tail
bound. It needs no per-neuron cancellation theorem.
\end{itemize}
